
Alignment in conversations is fundamentally bidirectional---both users and AI adapt to each other during multi-turn interactions. While current AI alignment methods focus on training models to match human preferences through RLHF, they do not measure whether genuine co-adaptation occurs during conversations. This work addresses this gap by proposing behavioral coupling as a metric for bidirectional alignment and providing empirical evidence that both components---user learning and AI responsiveness---exist in RLHF-trained systems.

Results confirm the theoretical foundation of bidirectional alignment. Users exhibit learning curves: reformulation rates decrease over turns (mean slope = -0.0214, p=0.012, Cohen's d=-0.246). AI systems exhibit behavioral responsiveness: response diversity correlates with query diversity (r=0.396, p<0.001). These findings, while showing smaller effect sizes than initially hypothesized, demonstrate that co-adaptation signals can be detected in real human-AI conversations using metrics extracted from conversation logs.

The temporal dynamics finding---responsiveness correlation strengthens from r=0.04 in short conversations to r=0.19 in longer conversations---provides evidence that alignment emerges through interaction rather than being a static property of AI policy quality. This pattern aligns with interactive alignment theory from psycholinguistics, where human conversation partners unconsciously synchronize linguistic representations through extended dialogue.

However, important limitations temper these findings. The small effect size for user learning (d=-0.246) and the median slope = 0 heterogeneity pattern suggest learning is conditional rather than universal. The correlation threshold failure for AI responsiveness (r=0.396 < 0.4) indicates responsiveness is present but weaker than predicted, potentially due to lexical diversity metrics that may underestimate semantic responsiveness. Most critically, the core coupling hypothesis---that correlation between learning rate and responsiveness predicts conversation quality beyond individual metrics---remains untested.

Despite these limitations, this work makes three contributions to AI alignment research. First, it introduces an operationalization of bidirectional alignment through observable conversation behaviors rather than subjective ratings or aggregate task metrics. Second, it demonstrates these metrics can be extracted from conversation logs without additional human evaluation, enabling scalable monitoring of alignment quality in deployed systems. Third, it provides empirical evidence that both user adaptation and AI responsiveness occur in RLHF-trained systems, validating the bidirectional framework's theoretical foundation and motivating future research on coupling effects.

Alignment is not something trained into AI systems---it is something that emerges during conversations when both user and AI adapt to each other. By measuring and optimizing for this co-adaptation, the field can move beyond asking whether AI produces responses humans prefer and toward asking whether human-AI interactions enable both participants to communicate effectively. That is the promise of bidirectional alignment.
